% !TeX root = ../navier-stokes-manuscript.tex
\subsection{Why the finite pieces are rectangles}\label{sub:BodyFiniteAdjacentRectangles}

The missing reach is a trace problem.
Fix one response \(V_n\) and one spatial height \(M\).
To carry \(V_n\) through the terminal interval in \(H^M\), its middle norm must be controlled by two neighboring coordinates of the grid: \(V_n\) one spatial step above it and \(V_{n+1}=\partial_tV_n\) one spatial step below it.

Let \(A=(1-\Delta)^{1/2}\).
For a strict cutoff \(T<T_*\), a smooth Hilbert-valued function \(v\), and \(0\leq a<b\leq T\),
\begin{equation}\label{eq:BodyAdjacentTraceIdentity}
  \norm{v(b)}_{H^M}^2-
  \norm{v(a)}_{H^M}^2
  =2\operatorname{Re}
  \int_a^b
  \pair{A^{M-1}v_t(t)}{A^{M+1}v(t)}_{L^2}\,dt.
\end{equation}
Cauchy--Schwarz places the change of the middle norm between those two coordinates:
\[
  \left|
  \norm{v(b)}_{H^M}^2-
  \norm{v(a)}_{H^M}^2
  \right|
  \leq
  2\norm{v_t}_{L^2(a,b;H^{M-1})}
   \norm{v}_{L^2(a,b;H^{M+1})}.
\]
Thus
\begin{equation}\label{eq:BodyAdjacentTracePair}
  v\in L^2(0,T;H^{M+1}),
  \qquad
  v_t\in L^2(0,T;H^{M-1})
\end{equation}
gives
\begin{equation}\label{eq:BodyAdjacentTraceConclusion}
  v\in C([0,T];H^M).
\end{equation}
Time mollification and Fourier truncation justify
\eqref{eq:BodyAdjacentTraceIdentity} for the pair in
\eqref{eq:BodyAdjacentTracePair}; weak \(H^M\) continuity together with continuity
of the norm gives the strong representative.

At \(M=2\), begin with \(V_0=u\):
\begin{equation}\label{eq:FirstConcreteAdjacentPair}
  u\in L^2(0,T;H^3),
  \qquad
  u_t\in L^2(0,T;H^1)
  \quad\Longrightarrow\quad
  u\in C([0,T];H^2).
\end{equation}
Move down to \(V_1=u_t\).
The same rule pairs \(V_1\in L^2H^3\) with \(V_2\in L^2H^1\) and carries \(V_1\) through \(H^2\).
The row \(V_1\) is shared: it is the lower input for the trace of \(V_0\), then the upper input for its own trace.

\Needspace{12\baselineskip}
\begin{center}
\captionsetup{hypcap=false}
\captionof{table}{The first two adjacent trace pairs at spatial height \(M=2\).}
\label{tab:FirstAdjacentRectangle}
\renewcommand{\arraystretch}{1.35}
\begin{tabular}{@{}c c c c@{}}
\toprule
response traced & one step above & next response below & middle trace \\
\midrule
\(V_0=u\) & \(V_0\in L^2H^3\) & \(V_1\in L^2H^1\) & \(V_0\in CH^2\) \\
\(V_1=u_t\) & \(V_1\in L^2H^3\) & \(V_2\in L^2H^1\) & \(V_1\in CH^2\) \\
\bottomrule
\end{tabular}
\end{center}

Every new temporal response repeats this overlap.
Stack the pairs for \(0\leq n\leq N\).
The squared velocity norm of the resulting finite rectangle is
\begin{equation}\label{eq:BodyRectangleNorm}
\begin{aligned}
  \mathfrak R_{N,M}(u;T)
  :={}&
  \sum_{n=0}^{N}
  \norm{V_n}_{L^2(0,T;H^{M+1})}^2
  \\
  &+
  \sum_{n=0}^{N}
  \norm{V_{n+1}}_{L^2(0,T;H^{M-1})}^2.
\end{aligned}
\end{equation}
Temporal depth \(N\) tells how many middle traces are selected.
Spatial height \(M\) tells where the same trace stencil sits in every selected row.
The two choices remain independent; when \(M=0\), the lower edge is read in \(H^{-1}\).

The pressure rows belong to the same finite history.
Write \(\cR_{N,M}[u_0;T]\) for the velocity coordinates measured by
\eqref{eq:BodyRectangleNorm}, together with the corresponding \(Q_n\)-coordinates reconstructed by
\eqref{eq:BodyMixedPressureRecurrence}.

For every strict cutoff, classical smoothness makes each fixed rectangle finite:
\begin{equation}\label{eq:StrictSubintervalRectangleBound}
  \mathfrak R_{N,M}(u;T)<\infty
  \qquad\text{for finite }N,M\text{ and }T<T_*.
\end{equation}
This still allows the rectangle to grow without bound as \(T\uparrow T_*\):
\[
  \forall\,T<T_*:\quad \mathfrak R_{N,M}(u;T)<\infty
  \quad\not\Longrightarrow\quad
  \sup_{0<T<T_*}\mathfrak R_{N,M}(u;T)<\infty.
\]
The trace calculation tells us what the second statement would buy.
For each fixed \(N\) and \(M\), it would carry every selected \(V_n\) through \(H^M\) on the whole terminal interval.
It does not produce that bound.
Neither the pressure recurrence nor overlap compatibility produces it.
The remaining step is to derive the cutoff-uniform estimate from the fixed datum \((u_0,\nu)\).
Everything that follows from that estimate remains conditional until this derivation is supplied.
